\documentclass[lualatex,ja=standard,a4paper,12pt]{bxjsarticle}
\usepackage{amsmath, amssymb, amsthm}
\usepackage{geometry}
\usepackage{xcolor}
\usepackage{tcolorbox}
\usepackage{booktabs}
\usepackage{array}

\geometry{margin=25mm}

\tcbuselibrary{theorems, skins, breakable}

\newtcolorbox{defbox}[1]{
  colback=blue!5, colframe=blue!60!black,
  fonttitle=\bfseries, title={#1}, arc=3pt, breakable
}
\newtcolorbox{thmbox}[1]{
  colback=green!5, colframe=green!50!black,
  fonttitle=\bfseries, title={#1}, arc=3pt, breakable
}
\newtcolorbox{pointbox}{
  colback=orange!8, colframe=orange!70!black,
  fonttitle=\bfseries, title={ポイント}, arc=3pt, breakable
}
\newtcolorbox{exbox}[1]{
  colback=gray!6, colframe=gray!50,
  fonttitle=\bfseries, title={例：#1}, arc=3pt, breakable
}

\title{補足資料 \\[6pt]
       \textbf{十分統計量（Sufficient Statistic）} \\[4pt]
       {\normalsize ゼミ資料}}
\author{}
\date{}

\begin{document}
\maketitle
\thispagestyle{empty}

% ===========================
\section{直感的な動機}
% ===========================

データ $y_1, y_2, y_3 = 3, 5, 7$ から母平均 $\theta$ を推定したいとする．

このとき，3つの数字を全部記録しておく必要があるだろうか？

\medskip
正規モデル $y_i \sim \mathrm{Normal}(\theta, \sigma^2)$ では，
$\theta$ に関する情報はすべて\textbf{標本平均 $\bar{y} = 5$} に集約されている．
データの順番や個々の値は $\theta$ の推定に追加の情報をもたらさない．

\medskip
この「情報の集約」を正確に表現したのが\textbf{十分統計量}の概念である．

% ===========================
\section{一般的な定義}
% ===========================

\begin{defbox}{十分統計量（Sufficient Statistic）}
統計量 $T(y_1,\ldots,y_n)$ が $\theta$ の\textbf{十分統計量}であるとは，

\[
  p(y_1,\ldots,y_n \mid T(y_1,\ldots,y_n) = t,\; \theta)
  \;\text{が}\; \theta \;\text{に依存しない}
\]

ことをいう．つまり，$T$ の値を知った後では，
生データを追加で見ても $\theta$ についての情報が増えない．
\end{defbox}

\medskip
\noindent この定義は直感と対応している：

\begin{itemize}
  \item $T$ に $\theta$ の情報がすべて入っている
        $\;\Rightarrow\;$ $T$ を固定すると，残りのデータのバラつきは $\theta$ と無関係
        $\;\Rightarrow\;$ 条件付き分布が $\theta$ に依存しない
\end{itemize}

% ===========================
\section{因数分解定理（確認に使う定理）}
% ===========================

定義を直接確認するのは難しいため，実用的には次の定理を使う．

\begin{thmbox}{Fisher--Neyman の因数分解定理}
$T(y_1,\ldots,y_n)$ が $\theta$ の十分統計量である
$\;\Longleftrightarrow\;$ 同時密度（または確率質量関数）が

\[
  p(y_1,\ldots,y_n \mid \theta)
  = g\!\bigl(T(y_1,\ldots,y_n),\; \theta\bigr) \times h(y_1,\ldots,y_n)
\]

と因数分解できる．ここで $g$ は $T$ と $\theta$ のみに依存し，
$h$ は $\theta$ に依存しない．
\end{thmbox}

\begin{pointbox}
確認の手順：同時密度を書き下して，
「$\theta$ を含む部分」と「$\theta$ を含まない部分」に分けられるか確かめる．
$\theta$ を含む部分がデータに依存する形が $T$ になっている．
\end{pointbox}

\subsection*{なぜ因数分解定理が定義と同値なのか}

同値を示すには2方向を示す必要がある．

\paragraph{方向1：因数分解できる $\Rightarrow$ 定義を満たす}

$p(y_1,\ldots,y_n|\theta) = g(T,\theta) \times h(y)$ と分解できるとして，
条件付き分布を計算する：

\[
  p(y \mid T=t, \theta)
  = \frac{p(y,\, T=t \mid \theta)}{p(T=t \mid \theta)}.
\]

分子は $T(y)=t$ を満たすとき $p(y|\theta) = g(t,\theta) \times h(y)$．

分母は $T=t$ になるデータを全部足し合わせたものなので，

\[
  p(T=t \mid \theta)
  = \sum_{\{y:\,T(y)=t\}} g(t,\theta) \times h(y)
  = g(t,\theta) \times \sum_{\{y:\,T(y)=t\}} h(y).
\]

割り算すると $g(t,\theta)$ が約分されて，

\[
  p(y \mid T=t, \theta)
  = \frac{g(t,\theta) \times h(y)}{g(t,\theta) \times \displaystyle\sum_{\{y:\,T(y)=t\}} h(y)}
  = \frac{h(y)}{\displaystyle\sum_{\{y:\,T(y)=t\}} h(y)}.
\]

\textbf{$\theta$ が消えた}．これがまさに定義の内容である．

\begin{pointbox}
核心は「$g(t,\theta)$ が分子と分母で同じ値なので約分できる」こと．\\
$\theta$ の情報は全部 $T$ に乗っており，$T$ を固定したら $\theta$ が消える
という構造が，計算の上で直接見えている．
\end{pointbox}

\paragraph{方向2：定義を満たす $\Rightarrow$ 因数分解できる}

$T$ が十分統計量（定義を満たす）とする．
すなわち $p(y \mid T=t, \theta) = q(y \mid t)$（$\theta$ に依存しない）と書ける．

尤度を次のように分解する：

\[
  p(y \mid \theta)
  = p(y \mid T=T(y),\, \theta) \times p(T=T(y) \mid \theta).
\]

第1因子は定義より $q(y \mid T(y))$（$\theta$ に依存しない）なので，

\[
  h(y) = q(y \mid T(y)),
  \qquad
  g(T(y),\, \theta) = p(T=T(y) \mid \theta)
\]

とおけば，

\[
  p(y \mid \theta) = g(T(y),\, \theta) \times h(y)
\]

という因数分解が得られる．\qed

% ===========================
\section{例で確認する}
% ===========================

\begin{exbox}{二項モデル：$T = \sum y_i$}

$y_1,\ldots,y_n \overset{\text{i.i.d.}}{\sim} \mathrm{Bernoulli}(\theta)$ のとき，
同時確率質量関数は

\[
  p(y_1,\ldots,y_n|\theta)
  = \prod_{i=1}^n \theta^{y_i}(1-\theta)^{1-y_i}
  = \theta^{\sum y_i}(1-\theta)^{n - \sum y_i}.
\]

因数分解すると

\[
  = \underbrace{\theta^{\sum y_i}(1-\theta)^{n-\sum y_i}}_{g(T,\,\theta),\; T=\sum y_i}
  \;\times\;
  \underbrace{1}_{h(y_1,\ldots,y_n)}.
\]

$h$ が定数なので，$T = \sum y_i$ が十分統計量である．

\medskip
\noindent\textbf{直感：}「1が何回出たか」さえ知れば，
どの順番で出たかは $\theta$ の推定に関係ない．
\end{exbox}

\begin{exbox}{ポアソンモデル：$T = \sum y_i$}

$y_1,\ldots,y_n \overset{\text{i.i.d.}}{\sim} \mathrm{Poisson}(\theta)$ のとき，

\[
  p(y_1,\ldots,y_n|\theta)
  = \prod_{i=1}^n \frac{\theta^{y_i} e^{-\theta}}{y_i!}
  = \underbrace{e^{-n\theta}\,\theta^{\sum y_i}}_{g(T,\,\theta)}
    \;\times\;
    \underbrace{\prod_{i=1}^n \frac{1}{y_i!}}_{h(y_1,\ldots,y_n)}.
\]

よって $T = \sum y_i$ が十分統計量．
\end{exbox}

\begin{exbox}{正規モデル：$T = (\bar{y},\, s^2)$}

$y_1,\ldots,y_n \overset{\text{i.i.d.}}{\sim} \mathrm{Normal}(\theta, \sigma^2)$（$\sigma^2$ 既知）のとき，

\[
  p(y_1,\ldots,y_n|\theta)
  \propto \exp\!\left\{-\frac{\sum(y_i - \theta)^2}{2\sigma^2}\right\}.
\]

指数部分を展開すると

\[
  \sum(y_i - \theta)^2
  = \sum y_i^2 - 2\theta \sum y_i + n\theta^2,
\]

なので

\[
  p(y_1,\ldots,y_n|\theta)
  = \underbrace{\exp\!\left\{\frac{\theta \sum y_i}{\sigma^2}
    - \frac{n\theta^2}{2\sigma^2}\right\}}_{g(\sum y_i,\,\theta)}
    \;\times\;
    \underbrace{\exp\!\left\{-\frac{\sum y_i^2}{2\sigma^2}\right\}}_{h(y_1,\ldots,y_n)}.
\]

よって $T = \sum y_i$（あるいは同値な $\bar{y}$）が十分統計量．

$\sigma^2$ も未知のときは $T = (\sum y_i,\, \sum y_i^2)$，
すなわち $(\bar{y},\, s^2)$ が2次元の十分統計量になる．

\end{exbox}

% ===========================
\section{指数型分布族との関係}
% ===========================

3.3節で学んだ指数型分布族 $p(y|\phi) = h(y)c(\phi)e^{\phi\,t(y)}$ において，
因数分解定理を適用すると：

\[
  p(y_1,\ldots,y_n|\phi)
  = \underbrace{c(\phi)^n \exp\!\left\{\phi \sum t(y_i)\right\}}_{g\left(\sum t(y_i),\,\phi\right)}
    \;\times\;
    \underbrace{\prod_{i=1}^n h(y_i)}_{h(y_1,\ldots,y_n)}.
\]

\begin{thmbox}{指数型分布族における十分統計量}
1次元指数型分布族では，$T = \sum_{i=1}^n t(y_i)$ が $\phi$ の十分統計量である．
\end{thmbox}

これは3つのモデルに一致する：

\medskip
\begin{center}
\renewcommand{\arraystretch}{1.5}
\begin{tabular}{lccc}
\toprule
モデル & $t(y)$ & 十分統計量 $T$ & 同値な統計量 \\
\midrule
二項 & $y$ & $\sum y_i$ & （1の個数） \\
ポアソン & $y$ & $\sum y_i$ & $\bar{y}$（標本平均） \\
正規（$\sigma^2$既知） & $y$ & $\sum y_i$ & $\bar{y}$ \\
正規（$\sigma^2$未知） & $(y,\, y^2)$ & $(\sum y_i,\, \sum y_i^2)$ & $(\bar{y},\, s^2)$ \\
\bottomrule
\end{tabular}
\end{center}

% ===========================
\section{ベイズ推論との関係}
% ===========================

十分統計量が重要なのは，ベイズ推論でも事後分布が十分統計量のみに依存するからである：

\[
  p(\phi \mid y_1,\ldots,y_n)
  \propto p(\phi) \times p(y_1,\ldots,y_n|\phi)
  = p(\phi) \times g(T, \phi) \times h(y_1,\ldots,y_n).
\]

$h$ は $\phi$ に依存しないので，

\[
  p(\phi \mid y_1,\ldots,y_n) \propto p(\phi) \times g(T, \phi).
\]

\begin{pointbox}
事後分布は生データ $(y_1,\ldots,y_n)$ ではなく，
\textbf{十分統計量 $T$ だけに依存する}．

つまり，ベイズ更新において「情報の圧縮」が保証されている．
共役事前分布のハイパーパラメータの更新則が $T$ だけを使う
シンプルな形になるのも，この構造が根拠になっている．
\end{pointbox}

% ===========================
\section{まとめ}
% ===========================

\begin{thmbox}{十分統計量のまとめ}
\begin{enumerate}
  \item \textbf{定義}：$T$ を知った後に生データを見ても $\theta$ の情報が増えない
        $\;\Leftrightarrow\;$ $T$ を条件付けた生データの分布が $\theta$ に依存しない．

  \item \textbf{確認方法}：因数分解定理 --- 同時密度が
        $g(T,\theta) \times h(y)$ と分解できれば $T$ が十分統計量．

  \item \textbf{指数型分布族}：$p(y|\phi) = h(y)c(\phi)e^{\phi t(y)}$ では
        $T = \sum t(y_i)$ が自動的に十分統計量になる．

  \item \textbf{ベイズ推論}：事後分布は十分統計量のみに依存するので，
        データを $T$ に圧縮してから更新すれば十分である．
\end{enumerate}
\end{thmbox}

\end{document}
